\documentclass[11pt]{article}
\usepackage[margin=1.05in]{geometry}
\usepackage{amsmath,amssymb,amsthm}
\newtheorem{theorem}{Theorem}
\newtheorem{lemma}[theorem]{Lemma}
\theoremstyle{definition}
\newtheorem{definition}[theorem]{Definition}
\newtheorem{example}[theorem]{Example}
\title{Temporal graphs, from snapshots to optimal paths}
\author{Michelle Döring \and Hendrik Molter}
\date{}
\begin{document}
\maketitle

\section{One graph, three representations}

% lax begin Lax330177.Representations
\begin{definition}
Fix a finite vertex set and a lifetime $T$. A temporal graph is a sequence
of loopless snapshots indexed by $0,\ldots,T-1$. In the directed setting
each snapshot records arcs. In the undirected setting the two orientations
represent the same physical edge. The \emph{footprint} contains the edges
that appear at least once. Give each such edge the set of times at which it
appears. Alternatively, list the arc occurrences $(t,u,v)$ in chronological
order, with ties in any order. All three descriptions retain the vertex set
and lifetime, including isolated vertices and empty final snapshots.
\end{definition}
% lax end

% lax begin Lax330177.SnapshotsLabels
\begin{lemma}
Converting snapshots to a labelled footprint and back recovers the original
temporal graph.
\end{lemma}
% lax end

% lax begin Lax330177Proofs.Basic.snapshots_labels
\begin{proof}
Fix a time $t$ and vertices $u,v$. Reconstructing the snapshot puts the arc
$(u,v)$ at $t$ exactly when its label set contains $t$. That is precisely
the condition for including $t$ when the label set was formed. Thus every
snapshot is recovered, and the lifetime was copied by both conversions.
\end{proof}
% lax end

% lax begin Lax330177.Settings
\begin{definition}
A labelling is \emph{simple} if each footprint edge has exactly one time,
\emph{proper} if arcs $(u,v)$ and $(v,w)$ with $w\ne u$ always have distinct labels,
and \emph{happy} if both hold. Opposite directed arcs may appear simultaneously. In symmetric snapshots the condition says that each snapshot is a matching.
\end{definition}
% lax end

% lax begin Lax330177.Paths
\begin{definition}
A temporal walk follows edges at their appearance times. The departure
times increase strictly in a \emph{strict} walk and weakly in a
\emph{nonstrict} one. A path is a walk whose vertices are pairwise distinct.
A vertex reaches another if a path joins them in the chosen mode.
A path with no edges reaches its own starting vertex.
\end{definition}
% lax end

% lax begin Lax330177.ReachabilityGraph
\begin{definition}
The \emph{reachability graph} in a chosen mode is the static digraph on the
same vertices with an arc $u\to v$ exactly when $u\ne v$ and a temporal path
in that mode joins $u$ to $v$. The zero-edge path gives reflexive reachability,
but the digraph omits loops.
\end{definition}
% lax end
\newpage
\section{When do strict and nonstrict paths agree?}

% lax begin Lax330177.ProperPathTimes
\begin{lemma}
On a proper temporal graph, every nonstrict path already has strictly
increasing departure times.
\end{lemma}
% lax end

% lax begin Lax330177Proofs.Basic.proper_path_times
\begin{proof}
Consecutive edges share their middle vertex. Their outer vertices are distinct because a path never repeats a vertex. Properness therefore
rules out equal departure times on consecutive edges. Weak increase at
every step becomes strict increase.
\end{proof}
% lax end

% lax begin Lax330177.ProperModes
\begin{theorem}
On a proper directed or undirected temporal graph, $u$ reaches $v$ strictly
if and only if $u$ reaches $v$ nonstrictly.
\end{theorem}
% lax end

% lax begin Lax330177Proofs.Basic.proper_modes
\begin{proof}
A strict path satisfies the nonstrict rule. In the other direction, the
preceding lemma shows that the very same nonstrict path, with its original
departure times, is a strict witness.
\end{proof}
% lax end

% lax begin Lax330177.BounceSameTime
\begin{example}
The corresponding assertion about every nonstrict \emph{walk} fails even
for a single undirected edge.
\end{example}
% lax end

% lax begin Lax330177Proofs.Basic.bounce_same_time
\begin{proof}
Let the edge appear at times 0 and 1. The graph is proper because it has
only one physical edge. At time 0 a walk can cross that edge and return
immediately. The two departures have the same time.
\end{proof}
% lax end

% lax begin Lax330177.Durations
\begin{definition}
Give each edge occurrence a nonnegative integer traversal time. A departure
at $a$ with traversal time $c$ arrives at $a+c$. A subsequent step may
depart upon arrival or after waiting. The lifetime bounds departures, and
the final arrival may follow the last snapshot.
\end{definition}
% lax end

% lax begin Lax330177.UniformPaths
\begin{lemma}
For fixed vertices and departure times, uniform traversal time zero gives
exactly the nonstrict paths, and uniform time one gives exactly the strict paths.
\end{lemma}
% lax end

% lax begin Lax330177Proofs.Basic.uniform_paths
\begin{proof}
Availability and vertex distinctness are the same in both descriptions.
For departures $a,b$ in traversal order, feasibility means $a+c\le b$.
For $c=0$ this says $a\le b$. For $c=1$ it says $a<b$, since the times
are integers.
\end{proof}
% lax end

\newpage
\section{What makes a path optimal?}

% lax begin Lax330177.OptimalPaths
\begin{definition}
Fix endpoints, a release time and an arrival deadline. Among feasible paths,
a \emph{shortest} path minimizes hops, a \emph{fastest} path minimizes
elapsed time, and a \emph{foremost} path minimizes arrival time.
A \emph{latest-departure} path maximizes its first departure time.
A \emph{minimum-waiting} path minimizes the sum of waits between successive
traversals. Waiting before the first departure contributes nothing.
\end{definition}
% lax end

% lax begin Lax330177.ElapsedEq
\begin{theorem}
If each traversal takes $c$ units, a path with $k$ hops and total internal
waiting time $w$ has elapsed time $w+kc$.
\end{theorem}
% lax end

% lax begin Lax330177Proofs.UniformCosts.elapsed_eq
\begin{proof}
Partition the interval from first departure to final arrival at the
successive arrivals and departures. The traversal intervals contribute
$kc$ in total. The remaining gaps are exactly the internal waits, of total
length $w$. On a zero-hop path departure and arrival coincide and both
terms vanish.
\end{proof}
% lax end

% lax begin Lax330177.ZeroOptima
\begin{lemma}
At uniform traversal time zero, a feasible path is fastest exactly when it
has minimum total waiting time.
\end{lemma}
% lax end

% lax begin Lax330177Proofs.UniformCosts.zero_optima
\begin{proof}
Set $c=0$ in the preceding identity. Every feasible path then has the same
value under both objectives, so the two sets of minimizers agree.
\end{proof}
% lax end

% lax begin Lax330177.NoWaitDetour
\begin{example}
Take a directed triangle with arcs $0\to1$, $1\to2$, $0\to2$, each
available at times 0 through 3. Let each traversal take $c\in\{0,1\}$
units, with release time 0 and deadline $3+c$. The two-hop detour
departing at $0,c$ minimizes waiting, uses more hops than the direct arc,
and is fastest exactly when $c=0$.
\end{example}
% lax end

% lax begin Lax330177Proofs.CostExamples.no_wait_detour
\begin{proof}
The second departure coincides with the first arrival, so the detour waits
zero time. Its two traversals take $2c$ units in total. At $c=0$ this is
the smallest possible elapsed time. At $c=1$ the direct arc needs only
one unit and arrives sooner when both paths start at time 0.
\end{proof}
% lax end
\end{document}
